\documentclass[11pt,a4paper]{article}
\usepackage[top=1.7cm,bottom=1.5cm,left=1.5cm,right=1.5cm,headsep=8pt,footskip=16pt]{geometry}
\usepackage{enumitem}
\usepackage{amssymb}
\usepackage{array}
\usepackage[table]{xcolor}
\usepackage{fancyhdr}
\usepackage{lastpage}
\usepackage{multicol}
\usepackage{polyglossia}
\setdefaultlanguage[calendar=gregorian, numerals=maghrib]{arabic}
\setotherlanguage{english}
\IfFontExistsTF{Amiri}{%
  \newfontfamily\arabicfont[Script=Arabic, Scale=1.08]{Amiri}
}{%
  \newfontfamily\arabicfont[Script=Arabic, Scale=1.08, Path=../fonts/, Extension=.ttf,
    UprightFont=*-Regular, BoldFont=*-Bold, ItalicFont=*-Italic,
    BoldItalicFont=*-BoldItalic]{Amiri}
}
\IfFontExistsTF{AlBattar}{%
  \newfontfamily\khatt[Script=Arabic, Scale=1.15]{AlBattar}
}{%
  \newfontfamily\khatt[Script=Arabic, Scale=1.15, Path=../fonts/,
    BoldFont=ae_AlBattar.ttf]{ae_AlBattar.ttf}
}
\renewcommand{\headrulewidth}{0.8pt}
\renewcommand{\footrulewidth}{0.8pt}
\fancyhf{}
\rhead{\khatt المدرسة العليا للأساتذة ببوزريعة}
\lhead{\khatt ملحقة بومرداس}
\lfoot{\khatt الأستاذ: د. بن عبد الرزاق يوسف}
\cfoot{\thepage/\pageref{LastPage}}
\pagestyle{fancy}
\parindent 0pt
\setlength{\columnsep}{0.6cm}
\setlength{\columnseprule}{0.4pt}
\setlist[enumerate,1]{label=\textbf{\arabic*.}, leftmargin=*, itemsep=5pt, topsep=2pt}
\begin{document}

\linespread{1.08}\selectfont
\begin{center}
{\LARGE\khatt\textbf{الاختبار الثاني في الذكاء الاصطناعي}}\\[4pt]
{\normalsize وضعيات إشكالية ــ أسئلة متعددة الاختيارات ــ الفصول من 1 إلى 6}
\end{center}
\vspace{-2pt}\rule{\textwidth}{0.6pt}\\[2pt]
\textbf{التعليمة:} يتضمن الاختبار خمس وضعيات إشكالية. اقرأ كل وضعية بعناية، ثم اختر التحليل أو التصرف الأنسب
من بين الاختيارات الخمسة؛ لكل وضعية إجابة صحيحة واحدة فقط.
اكتب \textbf{حرف} الإجابة الصحيحة في جدول الإجابات في آخر الاختبار.
لكل وضعية 4 نقاط (المجموع 20 نقطة).
\vspace{4pt}

\noindent
{\large\khatt\textbf{الوضعية 1}}\hfill{\small(الفصل 2)}\par\nopagebreak\vspace{3pt}
\setlength{\fboxsep}{7pt}\colorbox{gray!10}{\parbox{\dimexpr\linewidth-14pt}{لدى مهندس عشرون ألف نص من الكتب المدرسية غير مصنفة، ويريد بناء نموذج يصنّف كل نص إلى شعر عمودي أو شعر حر أو نثر، علمًا أن الشعر الحر قليل جدًا في هذه الكتب. اقترح زملاؤه خطوات مختلفة. ويريد نتائج يمكن الدفاع عنها. أي خطة تحترم منطق التعلم الموجَّه وتسمح بتقييم موثوق؟}}\par\nopagebreak
\begin{itemize}[leftmargin=1.0cm, labelsep=6pt, itemsep=3pt, topsep=4pt, parsep=0pt]
  \item[\textbf{أ)}] يصنّف يدويًا عينة كافية ومتوازنة من الأصناف الثلاثة، ويقسمها إلى جزء للتدريب وجزء للاختبار لا يراه النموذج أثناء التعلم، ثم يقيس الدقة على جزء الاختبار لكل صنف على حدة، لأن الدقة العامة قد تخفي ضعفًا في الصنف القليل.
  \item[\textbf{ب)}] يصنّف يدويًا عينة كافية من الأصناف الثلاثة، ويدرّب النموذج عليها كلها، ثم يقيس دقته على العينة نفسها، لأن استعمال كل الأمثلة المصنفة في التدريب يمنح النموذج أكبر قدر من المعرفة، والقياس على البيانات نفسها يوفر الوقت عند تحليل نتائجه.
  \item[\textbf{ج)}] يكتب قواعد العروض والقافية في برنامج يصنّف النصوص آليًا، ثم يستعمل نتائجه تسمياتٍ لتدريب النموذج دون مراجعة بشرية، ويقيس الدقة بمقارنة النموذج بالبرنامج نفسه، لأن تطابقهما دليل على صحة التصنيف.
  \item[\textbf{د)}] يدرّب النموذج على النصوص العشرين ألفًا دون تصنيف لأن كثرة البيانات تعوّض غياب التسميات، ثم يصنّف بنفسه مئة نص للاختبار ويقيس الدقة العامة عليها، فإن تجاوزت تسعين بالمئة اعتمد النموذج لكل الأصناف.
  \item[\textbf{هـ)}] يصنّف يدويًا عينة يغلب عليها النثر لأنه الأكثر في الكتب، ويقسمها إلى تدريب واختبار، ويكتفي بالدقة العامة على جزء الاختبار، لأن ارتفاعها يدل على أن النموذج يميّز الأصناف الثلاثة بالكفاءة نفسها، فالدقة العامة مؤشر كافٍ في هذه الحالات.
\end{itemize}
\par\vspace{6pt}
\noindent
{\large\khatt\textbf{الوضعية 2}}\hfill{\small(الفصل 3)}\par\nopagebreak\vspace{3pt}
\setlength{\fboxsep}{7pt}\colorbox{gray!10}{\parbox{\dimexpr\linewidth-14pt}{طلبت أستاذة شرح الاستعارة لتلاميذ السنة الرابعة متوسط، فجاء الشرح دقيقًا علميًا لكنه طويل ومثقل بمصطلحات كالمكنية والتصريحية والترشيح، وأمثلته كلها من الشعر القديم. وهدفها في هذه الحصة أن يميّز التلاميذ الاستعارة من التشبيه فقط. وتريد شرحًا قابلًا للعرض في دقائق قليلة من بداية الحصة القادمة. ما المتابعة الأنسب في المحادثة نفسها؟}}\par\nopagebreak
\begin{itemize}[leftmargin=1.0cm, labelsep=6pt, itemsep=3pt, topsep=4pt, parsep=0pt]
  \item[\textbf{أ)}] «أعد الشرح في حدود مئة وخمسين كلمة، مركّزًا على الفرق بين التشبيه والاستعارة، مع ثلاثة أمثلة من الشعر الجاهلي الراقي»، لأن الأمثلة الرفيعة ترفع مستوى التلاميذ وتعوّدهم على الفصحى العالية منذ البداية.
  \item[\textbf{ب)}] «الشرح طويل ومعقد ولا يناسب تلاميذي، أعده بطريقة أفضل وأوضح وأقصر، وركّز على المهم فقط»، لأن الأداة تستطيع أن تستنتج من شرحها السابق ما يجب حذفه وما يجب الإبقاء عليه دون تحديد ذلك.
  \item[\textbf{ج)}] تفتح محادثة جديدة وتكتب «اشرح الاستعارة لتلميذ»، لأن الأداة تتأثر بإجابتها السابقة في المحادثة نفسها فتكرر الأسلوب المعقد مهما كانت التعليمات الجديدة واضحة، والطلب القصير يضمن شرحًا بسيطًا، وهي طريقة يوصي بها كثير من المستخدمين.
  \item[\textbf{د)}] «أعد الشرح بلغة أبسط تناسب تلاميذ الرابعة متوسط، واحتفظ بأنواع الاستعارة ومصطلحاتها لأنها من البرنامج، لكن استبدل أمثلة الشعر القديم بأمثلة من حياة التلميذ اليومية، ثم أضف تمرينًا تطبيقيًا قصيرًا.»
  \item[\textbf{هـ)}] «أعد الشرح في حدود مئة وخمسين كلمة، مركّزًا على الفرق بين التشبيه والاستعارة بحذف أحد الطرفين، دون مصطلحات أخرى، مع ثلاثة أمثلة من حياة التلميذ اليومية، يقابل كلَّ مثال منها تشبيهٌ يوضّح الفرق.»
\end{itemize}
\par\vspace{6pt}
\noindent
{\large\khatt\textbf{الوضعية 3}}\hfill{\small(الفصل 4)}\par\nopagebreak\vspace{3pt}
\setlength{\fboxsep}{7pt}\colorbox{gray!10}{\parbox{\dimexpr\linewidth-14pt}{تريد أستاذة الاستعانة بأداة ذكاء اصطناعي عامة مجانية لكتابة ملاحظات تقويمية فردية على إنشاءات ثلاثين تلميذًا. تتضمن الملفات الأسماء الكاملة، والعلامات، وملاحظات عن حالات اجتماعية وصحية لبعض التلاميذ. وهي تعلم أن الأداة تابعة لشركة أجنبية، وأن شروط استعمالها تسمح بتخزين المحادثات مدة غير محددة. ما الإجراء الأنسب لحماية المعطيات الشخصية مع الاستفادة من الأداة؟}}\par\nopagebreak
\begin{itemize}[leftmargin=1.0cm, labelsep=6pt, itemsep=3pt, topsep=4pt, parsep=0pt]
  \item[\textbf{أ)}] تمتنع عن استعمال الأداة نهائيًا، لأن أي نص يكتبه تلميذ معطًى شخصي لا تجوز معالجته بأي أداة رقمية، حتى بعد حذف الأسماء وكل المعلومات الدالة على الهوية، ولو أذنت المؤسسة بذلك صراحة.
  \item[\textbf{ب)}] تحذف الأسماء الكاملة فقط وتُبقي الملاحظات الاجتماعية والصحية، لأنها تساعد الأداة على تكييف الملاحظات مع ظروف كل تلميذ، ولا يمكن معرفة صاحب الإنشاء بعد حذف اسمه حتى لو بقيت تفاصيله الشخصية.
  \item[\textbf{ج)}] تعوّض الأسماء بالأحرف الأولى من الاسم واللقب وتُبقي باقي المعلومات، لأن الأحرف الأولى تكفي لإخفاء الهوية في كل الحالات، ثم تعتمد الملاحظات المولَّدة كما هي لأن الأداة اطلعت على ظروف التلاميذ كاملة.
  \item[\textbf{د)}] تعوّض الأسماء برموز تحتفظ بجدول مقابلتها عندها، وتحذف الملاحظات الاجتماعية والصحية التي لا تحتاجها المهمة، ثم ترفع نصوص الإنشاءات وحدها، وتراجع الملاحظات المولَّدة قبل إعادتها إلى التلاميذ بأسمائهم.
  \item[\textbf{هـ)}] ترفع الملفات كاملة بعد تفعيل خيار عدم استعمال المحادثات في التدريب، لأن هذا الخيار يجعل الأداة لا تحتفظ بأي معلومة، فتصبح الأسماء والملاحظات الصحية محمية تمامًا دون حاجة إلى حذف أي شيء منها.
\end{itemize}
\par\vspace{6pt}
\noindent
{\large\khatt\textbf{الوضعية 4}}\hfill{\small(الفصل 5)}\par\nopagebreak\vspace{3pt}
\setlength{\fboxsep}{7pt}\colorbox{gray!10}{\parbox{\dimexpr\linewidth-14pt}{أعدّ أستاذ درسًا عن الحكمة في شعر المتنبي، فاقترحت عليه الأداة خمسة أبيات منسوبة إليه. وجد ثلاثة منها في الديوان المحقق، ووجد الرابع في منتديات كثيرة منسوبًا إلى المتنبي دون إحالة، ولم يجد الخامس في أي مصدر. وسيحضر الدرس مفتش المادة. كيف يتصرف تصرفًا منهجيًا قبل عرض الأبيات؟}}\par\nopagebreak
\begin{itemize}[leftmargin=1.0cm, labelsep=6pt, itemsep=3pt, topsep=4pt, parsep=0pt]
  \item[\textbf{أ)}] يعتمد البيت الخامس بعد أن يطلب من الأداة شرحه، فإن جاء الشرح متسقًا مع حكمة المتنبي ثبتت نسبته، ويستبعد الرابع لأن المنتديات مصدر غير موثوق، ويكتفي بالثلاثة الموثقة مع الخامس، وبذلك يكون قد احتاط للأمر.
  \item[\textbf{ب)}] يعتمد الأبيات الثلاثة الموثقة، ويبحث عن الرابع في الطبعات المحققة والشروح، فإن لم يجده تجنّب نسبته ولو كثر تداوله، ويستبعد الخامس لأن غيابه عن كل المصادر يرجّح أنه مولَّد يحاكي أسلوب الشاعر.
  \item[\textbf{ج)}] يعتمد الأبيات الأربعة الأولى، لأن كثرة تداول البيت الرابع في المنتديات قرينة كافية على شهرته وصحة نسبته، ويستبعد الخامس وحده لأنه الوحيد الذي لم يظهر في أي مصدر مكتوب أو رقمي متاح.
  \item[\textbf{د)}] يعتمد الأبيات الخمسة كلها، لأن صحة ثلاثة أبيات من خمسة تدل على أن الأداة اعتمدت على الديوان، فيكون البيتان الباقيان من طبعة أخرى لم يطّلع عليها الأستاذ، والأداة أوسع اطلاعًا منه، فلا داعي لمزيد من البحث.
  \item[\textbf{هـ)}] يسأل الأداة عن مصدر البيتين الرابع والخامس، فإن ذكرت اسم كتاب ورقم صفحة اعتمدهما، لأن ذكر المصدر بدقة يدل على أن الأداة استرجعت المعلومة من مرجع حقيقي ولم تولّدها.
\end{itemize}
\par\vspace{6pt}
\noindent
{\large\khatt\textbf{الوضعية 5}}\hfill{\small(الفصل 6)}\par\nopagebreak\vspace{3pt}
\setlength{\fboxsep}{7pt}\colorbox{gray!10}{\parbox{\dimexpr\linewidth-14pt}{صوّرت باحثة عشر صفحات من مخطوط بخط مغربي، وطلبت من أداة متعددة الأنماط نسخها نصًا مكتوبًا. جاءت النسخة مقروءة، لكن المقارنة الأولية أظهرت تصحيفًا في بعض أسماء الأعلام، وكلمات تامة أضافتها الأداة في مواضع ممحوة من المخطوط. وتنوي نشر النص محققًا في مجلة علمية محكّمة. ما الخطة الأنسب لاستثمار هذه النسخة؟}}\par\nopagebreak
\begin{itemize}[leftmargin=1.0cm, labelsep=6pt, itemsep=3pt, topsep=4pt, parsep=0pt]
  \item[\textbf{أ)}] تعدّ النسخة مسودة، وتقابلها بالصور، وتصحّح الأعلام، لكنها تُبقي الكلمات المضافة دون إشارة، لأن الإشارة إلى المواضع الممحوة تُثقل النص المحقق، والقارئ لا يهتم بالفرق بين ما كُتب وما أُكمل.
  \item[\textbf{ب)}] تطلب من الأداة إعادة النسخ ثلاث مرات، وتعتمد في كل موضع القراءة التي اتفقت عليها المحاولات، لأن اتفاقها يُغني عن مقابلة الصور، ويكشف التصحيف والإضافات تلقائيًا دون جهد بشري.
  \item[\textbf{ج)}] تعدّ النسخة مسودة، وتقابلها بالصور كلمة كلمة، وتصحّح الأعلام من كتب التراجم، وتحذف الكلمات المضافة في المواضع الممحوة وتشير إلى الطمس بعلامة متفق عليها، لأن الأداة قد تكمل النص بالمحتمل لا بالمكتوب.
  \item[\textbf{د)}] تعتمد النسخة بعد تصحيح الأعلام فقط، وتحتفظ بالكلمات المضافة في المواضع الممحوة، لأن الأداة أكملتها بما يناسب السياق، والسياق دليل كافٍ على الكلمات الأصلية في المخطوطات الممحوة جزئيًا، كما يفعل كثير من المحققين.
  \item[\textbf{هـ)}] تتخلى عن النسخة وتنسخ المخطوط يدويًا من البداية، لأن وجود تصحيف في بعض الأعلام يعني أن الأداة لم تقرأ الخط المغربي أصلًا، وأن كل الكلمات المقروءة في النسخة مولَّدة لا علاقة لها بالصور.
\end{itemize}
\par\vspace{6pt}
\begin{center}
{\khatt\textbf{جدول الإجابات}}\\[3pt]
\renewcommand{\arraystretch}{1.6}
\begin{tabular}{|>{\centering\arraybackslash}m{1.9cm}|*{5}{>{\centering\arraybackslash}m{1.2cm}|}}
\hline
\textbf{رقم السؤال} & \textbf{1} & \textbf{2} & \textbf{3} & \textbf{4} & \textbf{5} \\ \hline
\textbf{الإجابة} &  &  &  &  &  \\[8pt] \hline
\end{tabular}

\vspace{8pt}
\renewcommand{\arraystretch}{2.2}
\begin{tabular}{|>{\centering\arraybackslash}m{7.5cm}|>{\centering\arraybackslash}m{3.2cm}|>{\centering\arraybackslash}m{3.2cm}|}
\hline
\textbf{الاسم واللقب} & \textbf{الفوج} & \textbf{العلامة} \\ \hline
 & & \hfill /20 \\ \hline
\end{tabular}
\end{center}
\end{document}
